\ifextended

\templateentry
{0-1 BFS}
{\texttt{bfs01(g,s)} para pesos \texttt{0/1} em \texttt{g[v] = \{to,w\}}; $O(V+E)$.}
{\texttt{bfs01(g,s)}: distâncias de \texttt{s}; \texttt{g[v]} contém \texttt{\{to,w\}}, $w\in\{0,1\}$. Inalcançáveis: $10^9$. $O(V+E)$.}
{graph/bfs01.cpp}

\templateentry
{Dijkstra}
{\texttt{dijkstra(g,s)} retorna \texttt{dist} para \texttt{g[v] = \{to,w\}} em $O((V+E)\log V)$.}
{\texttt{dijkstra(g,s)}: distâncias de \texttt{s}; \texttt{g[v]} contém \texttt{\{to,w\}}, $w\ge0$. Inalcançáveis: $2^{62}$. $O((V+E)\log V)$.}
{graph/dijkstra.cpp}

\templateentry
{Ciclo de peso médio mínimo}
{}
{\texttt{minMeanCycle(n,edges)} para arcos \texttt{(u,v,peso)}: retorna fração \texttt{\{numerador,\allowbreak denominador\}} não reduzida, denominador positivo, ou \texttt{nullopt} se acíclico. Aceita pesos negativos, laços, paralelas e grafo desconexo. $O(V(V+E))$ tempo e $O(V^2)$ memória.}
{graph/min-mean-cycle.cpp}

\fi

\templateentry
{Caminho Euleriano}
{\texttt{eulerWalk(g,m,s)}: \texttt{\{vértices,ids\}}, ou ambos vazios se impossível a partir de \texttt{s}. IDs em $[0,m)$; em grafo não direcionado, use o mesmo ID nas duas direções. $O(n+m)$.}
{\texttt{eulerWalk(g,m,s)}: \texttt{\{vértices,ids\}}, ou ambos vazios se impossível a partir de \texttt{s}. IDs em $[0,m)$; em grafo não direcionado, use o mesmo ID nas duas direções. $O(n+m)$.}
{graph/eulerian-walk.cpp}

\ifextended

\templateentry
{Componentes do grafo complementar}
{}
{\texttt{complementCC(g)}: listas de vértices das componentes do complemento de um grafo simples não direcionado. $O(n+m)$ tempo, $O(n)$ memória adicional.}
{graph/complement-components.cpp}

\fi

\templateentry
{Árvore de pontes}
{\texttt{comp[v] = }componente de \texttt{v} na árvore de pontes. Ponte \texttt{uv} sse \texttt{comp[u] != comp[v]}. $O(V+E)$.}
{\texttt{bridges(g)}: componentes após remover as pontes de um grafo não direcionado. Aresta $uv$ é ponte sse \texttt{comp[u]!=comp[v]}. $O(V+E)$.}
{graph/bridges.cpp}

\templateentry
{Árvore de bloco e corte}
{\texttt{bct(g)}: floresta bipartida entre vértices $[0,n)$ e blocos $[n,|t|)$. \texttt{|t[v]|>1} sse $v$ é articulação. Vértices isolados também geram um bloco. $O(V+E)$.}
{\texttt{bct(g)}: floresta bipartida entre vértices originais $[0,n)$ e blocos $[n,|t|)$. \texttt{|t[v]|>1} sse $v$ é articulação. Vértices isolados também geram um bloco. $O(V+E)$.}
{graph/block-cut-tree.cpp}

\templateentry
{Numeração st}
{Assume que o grafo é biconexo após adicionar a aresta st. $O(V + E)$.}
{\texttt{stNumber(g,s,t)}: ordem com $s$ primeiro, $t$ último e cada vértice interno com um vizinho anterior e outro posterior. Assume que o grafo é biconexo após adicionar a aresta \texttt{s-t}. É necessário considerar o caso com dois vértices e nenhuma aresta separadamente. $O(V + E)$.}
{graph/st-numbering.cpp}

\templateentry
{Componentes fortemente conexas e grafo de condensação}
{\texttt{scc(g)}: \texttt{\{cnt,comp\}}, número de SCCs e componente de cada vértice; $O(V+E)$. \texttt{condGraph(g,cnt,comp)}: DAG sem arestas repetidas; $O(V+E\log(E+1))$.}
{\texttt{scc(g)}: \texttt{\{cnt,comp\}}, número de SCCs e componente de cada vértice; $O(V+E)$. \texttt{condGraph(g,cnt,comp)}: DAG sem arestas repetidas; $O(V+E\log(E+1))$.}
{graph/scc.cpp}

\templateentry
{2-SAT}
{\texttt{TwoSat ts(n)}; \texttt{clause(a,va,b,vb)} adiciona $(x_a=\mathtt{va})\lor(x_b=\mathtt{vb})$, índices 0-based. \texttt{solve()} retorna viabilidade; só então leia \texttt{val}. $O(V+E)$.}
{\texttt{TwoSat ts(n)}; \texttt{clause(a,va,b,vb)} adiciona $(x_a=\mathtt{va})\lor(x_b=\mathtt{vb})$, índices 0-based. \texttt{solve()} retorna viabilidade; só então leia \texttt{val}. $O(V+E)$.}
{graph/2-sat.cpp}

\ifextended

\templateentry
{Floresta de grafo funcional}
{\texttt{par[i] = }imagem de \texttt{i}. \texttt{root[i] = }representante do ciclo/componente de \texttt{i}. $O(n)$.}
{\texttt{root[i]}: representante do ciclo/componente de \texttt{i}. $O(n)$.}
{graph/functional-forest.cpp}

\fi

\templateentry
{Dominator tree}
{\texttt{domTree(g,s)} retorna o dominador imediato; \texttt{-1} para inalcançáveis. Quase linear.}
{$u$ domina $v$ se todo caminho de $s$ a $v$ passa por $u$. \texttt{domTree(g,s)} retorna \texttt{idom[v]}, o dominador imediato de $v$; \texttt{idom[s]=s}, $-1$ para inalcançáveis. $O((V + E)\alpha(V))$ na prática.}
{graph/dominator-tree.cpp}

\templateentry
{Árvore geradora mínima direcionada}
{\texttt{dmst(n,r,e)} retorna o custo da arborescência mínima, ou \texttt{-1}; $O(E \log(V))$.}
{\texttt{dmst(n,r,e)}: custo da arborescência mínima enraizada em $r$, ou $-1$ se impossível. Arestas \texttt{E\{a,b,w\}}: $a\to b$ com custo $w$. $O(E \log(V))$.}
{graph/directed-mst.cpp}

\templateentry
{Matching bipartido máximo}
{\texttt{hopKarp(g,l,r)}: $g$ da esquerda para a direita; inicialize \texttt{l,r} com $-1$. Escreve parceiros (ou $-1$) e retorna o tamanho do matching. $O(E\sqrt V)$.}
{\texttt{hopKarp(g,l,r)}: $g$ da esquerda para a direita; inicialize \texttt{l,r} com $-1$. Escreve parceiros (ou $-1$) e retorna o tamanho do matching. $O(E\sqrt V)$.}
{graph/hopcroft-karp.cpp}

\templateentry
{Matching bipartido com menor custo}
{\texttt{hungarian(a)}: \texttt{\{custo,match\}}, uma coluna distinta por linha; matriz $n\times m$, $n\le m$. $O(n^2m)$.}
{\texttt{hungarian(a)}: \texttt{\{custo,match\}}, uma coluna distinta por linha; matriz $n\times m$, $n\le m$. $O(n^2m)$.}
{graph/hungarian.cpp}

\templateentry
{Encontrar matching geral}
{\texttt{matching(g)}: parceiros em grafo não direcionado, $-1$ se livre. $O(n^3)$.}
{\texttt{matching(g)}: parceiros em grafo não direcionado, $-1$ se livre. $O(n^3)$.}
{graph/matching.cpp}

\ifextended

\templateentry
{Matching geral}
{\texttt{matchingSize(n,es,its=2)}: tamanho estimado por matriz de Tutte; módulo primo ímpar. \textbf{Monte Carlo} com possível subestimação; \texttt{its} repetições. Ignora self-loops. $O(\mathtt{its}\cdot n^3)$ para módulo fixo.}
{\texttt{matchingSize(n,es,its=2)}: tamanho estimado por matriz de Tutte; módulo primo ímpar. \textbf{Monte Carlo} com possível subestimação; \texttt{its} repetições. Ignora self-loops. $O(\mathtt{its}\cdot n^3)$ para módulo fixo.}
{graph/matching-size.cpp}

\fi

\notebookentry{Fluxo máximo (Dinic)}
{\texttt{addEdge(u,v,cap)}: direcionada, capacidade não negativa.
\texttt{flow(s,t,mx)} retorna o fluxo adicional, até \texttt{mx}; requer $s\ne t$ e mantém o residual entre chamadas. $O(V^2E)$ para fluxo máximo.
\extendedonly{O fluxo de cada aresta original está em \texttt{e[id].f}. Após fluxo máximo sem limite, vértices alcançáveis por residual positivo formam o lado da fonte de um corte mínimo.}}
{\addtemplate{graph/dinic.cpp}}

\templateentry
{Árvore de Gomory--Hu}
{\texttt{gomoryHu(n,es)}: grafo não direcionado, arestas \texttt{(u,v,c)}, $c\ge0$. Retorna árvore onde o mínimo peso no caminho $u$--$v$ é seu corte mínimo. $n-1$ fluxos; desconexos ligados por peso zero.}
{Grafo não direcionado com capacidades não negativas; arestas \texttt{(u,v,c)}, paralelas permitidas. Retorna arestas da árvore: o menor peso no caminho entre $u,v$ é o valor do corte mínimo entre eles. $n-1$ fluxos máximos; vértices desconexos ficam ligados por peso zero.}
{graph/gomory-hu.cpp}

\notebookentry{Fluxo de custo mínimo}
{\texttt{addEdge(u,v,cap,cost)}; \texttt{flow(s,t,k)} retorna \texttt{\{fluxo,custo\}} adicional até $k$, mantendo o residual. Requer $s\ne t$, sem self-loops e sem ciclos negativos no residual inicial; custos negativos são permitidos. $O(VE+FE\log V)$, $F$ = número de aumentos.}
{\addtemplate{graph/mcf.cpp}}

\ifextended

\templateentry
{Grafo com sequência de graus}
{\texttt{havelHakimi(deg)}: arestas de um grafo simples com esses graus. Requer sequência gráfica válida; não detecta impossibilidade. $O(n+m)$.}
{\texttt{havelHakimi(deg)}: arestas de um grafo simples com esses graus. Requer sequência gráfica válida; não detecta impossibilidade. $O(n+m)$.}
{graph/havel-hakimi.cpp}

\templateentry
{Enumeração de triângulos}
{}
{\texttt{triangles(g,emit)}: grafo simples não direcionado. Chama \texttt{emit(u,v,w)} uma vez por triângulo; ordem dos vértices não garantida. $O(n+m\sqrt m)$ tempo para callback $O(1)$, $O(n+m)$ memória adicional.}
{graph/triangles.cpp}

\templateentry
{Coloração de arestas}
{}
{\texttt{edgeColoring(n,edges)}: cor de cada aresta na ordem de entrada, em $[0,\Delta]$, onde $\Delta$ é o grau máximo. Requer grafo não direcionado simples. Não garante o mínimo de cores. $O(nm+n)$, memória $O(n(\Delta+1))$.}
{graph/edge-coloring.cpp}

\fi
